\section{Research scope, conventions, and dependencies}
\label{sec:scope}

The repository has a substantial body of overlapping reports. A useful
continuation must first fix its comparison point. We use the commit on
the title page, whose last change added six new research bundles.
The canonical manuscript already incorporates many older results;
the six new bundles contain additional theorems that must also be
credited before identifying an open question. The baseline and its
mathematical status are recorded in \cite{Repo,Incoming} and in the
individual reports \cite{BTD,CTD,DHR,DPL,MR,PTI}.

\subsection{What is continued here}

The optimal-prefix question concerns the centered moments of two
ordinary zeta tails. The incoming reports prove the diagonal case,
the adjacent even-order case $p=q-2$, and the families with smaller
order $2,4,$ or $6$; they leave a general unequal-order determinant
unresolved \cite{BTD,DPL,PTI}. Section~\ref{gap:section} proves the
entire next gap $q-p=4$. Thus its first pair outside those earlier
small-order families is $(p,q)=(8,12)$. Section~\ref{tel:section}
then turns late-coordinate relations into explicit telescopers for
all even pairs, whether or not their optimal-prefix determinant is
known to be nonzero.

The inverse-Gamma normal form and the full digamma--trigamma sector
are inherited from \cite{DPL,PTI}. The explicit request in
\cite[research questions 4--5]{PTI} asks for the next mixed sector,
including arbitrary powers of the first two derivative variables.
Section~\ref{var:section} gives an exact recurrence for the entire
sector with largest derivative index at most two. It also proves
closed formulas for all terms containing exactly two or exactly three
copies of $\psi''$. This is a coefficient theorem in a differential
polynomial algebra; its integration consequences retain every endpoint
term.

The rank-two higher-resonance problem is stated explicitly in
\cite[Section 9]{DHR}. Section~\ref{res:section} treats that target,
including the logarithmic square in the hard cutoff and its difference
from the spectral Laurent constant. The logarithmically weighted
finite-part problem in \cite{DPL} is treated in
Section~\ref{log:section}. These are analytic continuations with
specified counterterms, so their normalization is part of each theorem.

\subsection{Ordinary tails, strict sums, and spectral derivatives}

For $s>1$ and $n\ge0$, put
\begin{equation}\label{scope:tails}
 H_n^{(s)}=\sum_{k=1}^n\frac1{k^s},\qquad
 A_s(n)=\zeta(s)-H_n^{(s)}=\zeta(s,n+1).
\end{equation}
For integer $s\ge2$ this gives
\[
 A_s(n)=\frac{(-1)^s}{(s-1)!}\psi^{(s-1)}(n+1).
\]
The even orders used in the tail sections therefore have positive
polygamma normalization. These are standard Hurwitz identities
\cite[Sections 5.15 and 25.11]{DLMF}.

A multiple zeta value uses strictly decreasing indices:
\[
 \zeta(s_1,\ldots,s_d)=
 \sum_{n_1>\cdots>n_d\ge1}
       n_1^{-s_1}\cdots n_d^{-s_d}
\]
where the series converges, with meromorphic continuation when needed.
A prime on a function of a spectral variable denotes differentiation
in that variable. In particular it never denotes differentiation of a
formula known only at integer index values.

We use generalized Stieltjes constants with the convention
\begin{equation}\label{scope:stieltjes}
 \zeta(1+u,a)=\frac1u+
       \sum_{j\ge0}\frac{(-1)^j}{j!}\gamma_j(a)u^j,
 \qquad a>0,
\end{equation}
and $\gamma_j=\gamma_j(1)$. Thus $\gamma_0(a)=-\psi(a)$.

\subsection{Two precise meanings of finite part}

For a function with a finite Laurent-logarithmic singular expansion at
zero, its original-coordinate integral finite part is
\begin{equation}\label{scope:fp}
 \FP\int_0^1f(x)\dd x
 =\CT_{\varepsilon\downarrow0}
             \int_\varepsilon^1f(x)\dd x.
\end{equation}
The symbol $\CT$ selects the constant multiplying
$\varepsilon^0(\log\varepsilon)^0$. The coordinate is $x$ itself.
For a polynomial $H$ in digamma derivatives this convention implies
\begin{equation}\label{scope:boundary}
 \FP\int_0^1\frac{d}{dx}H(\psi(x),\psi'(x),\ldots)\dd x
 =H(\psi(1),\psi'(1),\ldots)
   -[x^0]H(\psi(x),\psi'(x),\ldots).
\end{equation}
This follows directly by taking the constant in the endpoint difference.
A change of coordinate must be evaluated separately.

The centered tail moment is a different, explicitly specified object:
\begin{equation}\label{scope:centered}
 D_{p,q}(s)=\sum_{n\ge0}
       A_p(n)A_q(n)(n+\tfrac12)^{-s},
 \qquad \re s>3-p-q.
\end{equation}
We continue it meromorphically and write
$D_r=D_{p,q}(-2r)$. For positive even $p,q$, the centered asymptotic
expansion has only even powers; the possible poles are at odd
integers, so all these values are regular. Equivalently, subtract
from $(n+\tfrac12)^{2r}A_p(n)A_q(n)$ its polynomial part in
$n+\tfrac12$. The resulting summand is $O(n^{-2})$, and its ordinary
sum equals $D_r$ because
$\zeta(-2j,\tfrac12)=0$ for $j\ge0$. Section~\ref{tel:section}
proves this relation with the finite-cutoff constants made explicit.

\subsection{Status and reproducibility}

The proofs below establish their all-order claims algebraically or
analytically. Exact rational verification files check the matrix
conjugacies, polynomial certificates, normal forms, and displayed
specializations. Independent numerical computations compare different
analytic representations. They are diagnostics, not interval proofs,
unless a check is explicitly described as rational arithmetic.

No assertion of arithmetic independence follows from a formal-rank
statement. In particular, the symbols used as late zeta coordinates
are treated as formal basis labels even though their numerical values
have known multiplicative relations. Likewise, the new recurrence for
Gamma coefficients does not prove that the corresponding mixed moments
form a minimal numerical constant basis.
